\documentclass[12pt,a4paper]{article}

% ============================================================
% PACKAGES
% ============================================================
\usepackage[left=1in,right=1in,top=1in,bottom=1in]{geometry}
\usepackage{graphicx}
\usepackage{booktabs}
\usepackage{array}
\usepackage{tabularx}
\usepackage{colortbl}
\usepackage{xcolor}
\usepackage{makecell}
\usepackage{fancyhdr}
\usepackage{titlesec}
\usepackage{amsmath}
\usepackage{tcolorbox}

\definecolor{headerblue}{RGB}{213,220,228}

\newcolumntype{Y}{>{\centering\arraybackslash}X}
\newcolumntype{L}{>{\raggedright\arraybackslash}X}

\setlength{\parindent}{0pt}
\setlength{\parskip}{6pt}

\pagestyle{fancy}
\fancyhf{}
\rhead{SE: Machine Learning Lab}
\lhead{Lab 5}
\rfoot{\thepage}

\begin{document}

% ============================================================
% TABLE 1: DEPARTMENT HEADER
% ============================================================
\renewcommand{\arraystretch}{2.0}
\begin{tabularx}{\textwidth}{|Y|}
\hline
\rowcolor{headerblue} \textbf{\large Department of Computer and Software Engineering} \\
\hline
\textbf{\large SE: Machine Learning} \\
\hline
\end{tabularx}
\renewcommand{\arraystretch}{1.0}

\vspace{12pt}

% ============================================================
% TABLE 2: INSTRUCTOR INFO
% ============================================================
\renewcommand{\arraystretch}{2.0}
\begin{tabularx}{\textwidth}{|L|L|}
\hline
\textbf{Course Instructor:} Dr. Ahmed Raza & \textbf{Date:} \\
\hline
\textbf{Day:} & \textbf{Semester:} \\
\hline
\end{tabularx}
\renewcommand{\arraystretch}{1.0}

\vspace{12pt}

% ============================================================
% TITLE
% ============================================================
\begin{center}
{\Large\bfseries Lab 5: Decision Tree Split using Gini and Information Gain}
\end{center}

\vspace{6pt}

% ============================================================
% TABLE 3: LAB METADATA
% ============================================================
\renewcommand{\arraystretch}{1.8}
\begin{tabularx}{\textwidth}{|Y|c|c|c|c|}
\hline
\textbf{Lab Title} & \textbf{CLO} & \textbf{BT} & \textbf{PLO} & \textbf{Weightage} \\
\hline
\makecell{{\small Decision Tree Split} \\ {\small using Gini and Information Gain}} & & & & \\
\hline
\end{tabularx}
\renewcommand{\arraystretch}{1.0}

\vspace{6pt}

% ============================================================
% TABLE 4: STUDENT INFO
% ============================================================
\renewcommand{\arraystretch}{2.0}
\begin{tabularx}{\textwidth}{|Y|Y|Y|Y|Y|}
\hline
\textbf{Name} & \textbf{Roll Number} & \textbf{Report /10} & \textbf{Viva /5} & \textbf{Total /15} \\
\hline
 & & & & \\
\hline
\end{tabularx}
\renewcommand{\arraystretch}{1.0}

\vspace{12pt}

\begin{flushright}
Checked on: \_\_\_\_\_\_\_\_\_\_\_\_\_\_\_\_

\vspace{6pt}

Signature: \_\_\_\_\_\_\_\_\_\_\_\_\_\_\_\_
\end{flushright}

\newpage

% ============================================================
% LAB CONTENT
% ============================================================

\section{Learning Outcomes}

After completing this lab, the student will be able to:

\begin{itemize}
\item Compute Gini impurity for binary classification.
\item Calculate Information Gain from candidate splits.
\item Implement the split selection procedure used by decision trees.
\item Observe bias toward high-cardinality features.
\end{itemize}

\section{Dataset}

We will use the Iris dataset from \texttt{sklearn}.  
To simplify the problem we restrict the dataset to two classes.

Properties:

\begin{itemize}
\item Samples: 100
\item Features: 4
\item Classes: Binary
\end{itemize}

Dataset loading:

\begin{verbatim}
from sklearn.datasets import load_iris
\end{verbatim}

\section{Gini Impurity}

For binary classification:

\[
Gini = 1 - (p_0^2 + p_1^2)
\]

A perfectly pure node has Gini impurity equal to 0.

\section{Information Gain}

The reduction in impurity after a split is:

\[
Gain =
Gini(parent) -
\left(
\frac{n_L}{n}Gini(L)
+
\frac{n_R}{n}Gini(R)
\right)
\]

The split with the highest information gain is selected.

% ============================================================

\section{Part A: Implement Gini Impurity (3 Marks)}

\begin{tcolorbox}

Implement the function:

\begin{verbatim}
gini(y)
\end{verbatim}

The function should:

\begin{itemize}
\item compute class probabilities
\item apply the Gini formula
\item return impurity
\end{itemize}

Verify using:

\begin{verbatim}
y = [0,0,0,1]
\end{verbatim}

Compute impurity manually and compare.

\vspace{5cm}

\end{tcolorbox}

% ============================================================

\section{Part B: Implement Information Gain (4 Marks)}

\begin{tcolorbox}

Implement:

\begin{verbatim}
information_gain(parent, left, right)
\end{verbatim}

Steps:

\begin{itemize}
\item compute parent impurity
\item compute child impurities
\item compute weighted average
\item subtract from parent impurity
\end{itemize}

\vspace{5cm}

\end{tcolorbox}

% ============================================================

\section{Part C: Implement Tree Split Logic (4 Marks)}

\begin{tcolorbox}

Implement the function:

\begin{verbatim}
find_best_split(X, y)
\end{verbatim}

The function should:

\begin{itemize}
\item iterate over features
\item compute candidate thresholds
\item evaluate information gain
\item return the best feature and threshold
\end{itemize}

Use NumPy operations where possible.

\vspace{6cm}

\end{tcolorbox}

% ============================================================

\section{Part D: Cardinality Bias Experiment (4 Marks)}

\begin{tcolorbox}

A new feature will be added:

\begin{verbatim}
index_feature = sample_index / N
\end{verbatim}

This feature contains many unique values.

Run the split algorithm again and observe:

\begin{itemize}
\item which feature is selected
\item whether the high-cardinality feature is preferred
\end{itemize}

Explain why this occurs.

Explanation
1.Feature 2 is selected in both cases.
2.Whether the high-cardinality feature is preferred: No, it is not preferred.
3.While decision trees are theoretically biased toward high-cardinality features, this only manifests if the feature can isolate classes. Because train_test_split randomly shuffles the target y, the sequential high-cardinality index_feature cannot separate the classes with a single threshold split. Furthermore, Feature 2 provides a nearly perfect Information Gain (0.49875). In a single-node split, a perfectly predictive continuous feature will mathematically overpower a high-cardinality noise feature.

\vspace{6cm}

\end{tcolorbox}

\section{Expected Output}

Your program should print:

\begin{itemize}
\item Best split feature
\item Threshold
\item Information gain
\item Result of cardinality bias experiment
\end{itemize}

\newpage
% ============================================================
% Assessment Rubrics
% ============================================================
\begin{center}
{\Large\bfseries Assessment Rubrics}

\end{center}

\textbf{Method:} Lab reports and instructor observation during lab sessions

\textbf{Outcome Assessed:}

\begin{enumerate}
\item Ability to conduct experiments, as well as to analyze and interpret data (P).
\item Ability to use the techniques, skills, and modern engineering tools necessary for engineering practice (P).
\end{enumerate}

% ============================================================
% TABLE 5: Assessment Rubric Table
% ============================================================
{\scriptsize % <--- Sets font size smaller than normal text
\renewcommand{\arraystretch}{1.5}
\noindent
\begin{tabularx}{\textwidth}{|>{\raggedright\arraybackslash}p{0.14\textwidth}|>{\raggedright\arraybackslash}X|>{\raggedright\arraybackslash}X|>{\raggedright\arraybackslash}X|c|}
\hline
\textbf{Performance} & 
\textbf{Exceeds expectation (4-5)} & 
\textbf{Meets expectation (3-2)} & 
\textbf{Does not meet expectation (1)} & 
\textbf{Marks} \\
\hline

% Row 1
1. Realization of Experiment [a, b] & 
Selects relevant equipment to the experiment, develops setup diagrams of equipment connections or wiring. & 
Needs guidance to select relevant equipment and to develop equipment connection or wiring diagrams. & 
Incapable of selecting relevant equipment to conduct the experiment. & 
\\ \hline

% Row 2
2. Conducting Experiment [a, b] & 
Does proper calibration of equipment, carefully examines equipment moving parts, and ensures smooth operation. & 
Calibrates equipment, examines equipment moving parts, and operates the equipment with minor error. & 
Unable to calibrate appropriate equipment, and equipment operation is substantially wrong. & 
\\ \hline

% Row 3
3. Laboratory Safety Rules [a] & 
Respectfully and carefully observes safety rules and procedures. & 
Observes safety rules and procedures with minor deviation. & 
Disregards safety rules and procedures. & 
\\ \hline

% Row 4
6. Data Collection [a] & 
Plans data collection to achieve experimental objectives, and conducts an orderly and complete data collection. & 
Plans data collection to achieve experimental objectives, and collects complete data with minor error. & 
Does not know how to plan data collection; data collected is incomplete and contain errors. & 
\\ \hline

% Row 5
7. Data Analysis [a] & 
Accurately conducts simple computations and statistical analysis; correlates results to known theoretical values; accounts for measurement errors. & 
Conducts simple computations with minor error; reasonably correlates results to known theoretical values; attempts to account for errors. & 
Unable to conduct simple statistical analysis; no attempt to correlate or explain errors. & 
\\ \hline

% Row 6
8. Computer Use [a] & 
Uses computer to collect and analyze data effectively. & 
Uses computer to collect and analyze data with minor error. & 
Does not know how to use computer to collect and analyze data. & 
\\ \hline

% Total Row
\multicolumn{4}{|r|}{\textbf{Total}} & \\
\hline
\end{tabularx}
\renewcommand{\arraystretch}{1.0}
}

\textbf{Faculty:}

\textbf{Name:} \rule{5cm}{0.4pt}

\textbf{Signature:} \rule{5cm}{0.4pt}

\textbf{Date:} \rule{5cm}{0.4pt}

\end{document}